\subsection{Cross-Level Synthesis}
% Compare every level using Return, Persist, Gate, AI responsibility and authority, and remaining human responsibility. Summarize the capability gains and risks accompanying each transfer of responsibility, and distinguish levels with substantial public evidence from early prototypes, company disclosures, and research visions.

Across B0--L5, the autonomy hierarchy can be understood as a progressive transfer of responsibility over the improvement loop. The levels differ not primarily in the particular algorithm or artifact being optimized, but in which improvement decisions are internalized by the AI system, what state persists into later rounds, and which acceptance conditions remain externally protected.


The central boundaries between adjacent levels can therefore be stated compactly. The transition from B0 to L1 is \textbf{persistence}: an accepted change must survive the current task. L1 to L2 transfers \textbf{strategy selection}: AI decides which intervention to attempt rather than merely executing a prescribed one. L2 to L3 adds control over the \textbf{future learning agenda}, so the evolving learner influences what experience is acquired next. L3 to L4 extends this feedback process into \textbf{persistent deployment and environmental interaction}, where system changes must remain useful under changing operational conditions. L4 to L5 finally introduces \textbf{recursive inheritance}, in which the mechanism responsible for later improvement itself becomes an inherited target of improvement.

Importantly, a higher autonomy level does not by itself imply a better improvement process. Greater delegated authority can coexist with inefficient search, unreliable feedback, regression, evaluator exploitation, or poor transfer. We therefore separate the scope of improvement responsibility internalized by AI from the quality and efficiency of the resulting improvement. Across the evidence reviewed here, lower and intermediate levels are supported by a comparatively broad range of systems, whereas L3--L4 capabilities remain more domain-dependent and end-to-end L5 evidence is concentrated in bounded prototypes and emerging industrial or research systems. This distinction motivates the application-specific analysis in the following section.